%!TEX root =  tfg.tex
\chapter{Planificación}

\begin{quotation}[Novelist]{Ernest Hemingway (1899--1961)}
The good parts of a book may be only something a writer is lucky enough to overhear or it may be the wreck of his whole damn life -- and one is as good as the other.
\end{quotation}

\begin{abstract}
Este capítulo describe la estrategia de gestión y el cronograma seguidos 
para el desarrollo de \textit{Via Inferni}. El proyecto se organiza ahora en tres fases 
claras: planificación y gestión, diseño e implementación, y cierre del proyecto. 
La Fase 2 se desglosa además en varias iteraciones asociadas al backlog, de forma que 
quede trazabilidad entre cada issue, su motivación y el momento en el que se consideró 
realmente terminada. El objetivo es ofrecer una visión clara de cómo se han administrado 
los recursos y el tiempo para transformar los requisitos en un producto funcional.
\end{abstract}

\section{Resumen temporal del proyecto}

\begin{table*}[htb]
	\centering
	\begin{coolTable}{ll}{2}
{Resumen del proyecto}
	\textbf{Fecha de inicio}&10/10/2014\\
	\textbf{Fecha de fin}&10/10/2014\\
	\textbf{Periodicidad de las revisiones}&3 semanas\\
	\end{coolTable}
	\caption{Tabla resumen de tiempos y planificación}
\end{table*}

\section{Planificación inicial}

\textbf{ESTE CAPÍTULO DEBE ESCRIBIRSE AL COMIENZO DEL PROYECTO}

En esta sección se presenta el desglose inicialmente planificado de las fases en las que se ha dividido el proyecto,
así como una justificación de cada una de ellas. Se han definido tres fases principales:

\begin{table*}[htb]
	\centering
	\begin{coolTable}{ll}{2}
{Resumen de fases del proyecto}
	\textbf{Fase 1: Planificación y Gestión}&09/02/2026 a 15/02/2026\\
	\textbf{Fase 2: Diseño e Implementación}&16/02/2026 a 10/05/2026\\
	\textbf{Fase 3: Cierre del Proyecto}&11/05/2026 a 25/05/2026\\
	\end{coolTable}
	\caption{Planificación temporal de fases del proyecto}
\end{table*}

El trabajo se realizará del 9 de febrero al 25 de mayo, con una duración total de aproximadamente 
tres meses y medio, y un total de 99 días. Se estiman unas 3 o 4 horas de trabajo diario,
lo que se traduce en un total aproximado de 300 horas de trabajo por miembro. 

Cada integrante adaptará las horas de trabajo diarias a su disponibilidad, pero se mantendrá 
una comunicación constante con reuniones recurrentes para asegurar la integración. Al menos
una vez al final de cada semana se realizará una reunión de seguimiento para revisar el progreso y
planificar las tareas de la semana siguiente.

\subsection{Fase 1: Planificación y Gestión}

Esta fase consistirá en realizar un análisis detallado de los requisitos del proyecto, definir 
los objetivos específicos y elaborar un plan de trabajo detallado. Además, se fijará la estructura 
de gestión del proyecto, la organización del backlog, la cadencia de revisión y la forma en la que 
se documentarán los avances para que la memoria evolucione al mismo ritmo que el desarrollo.

Se ha estimado una duración de una semana para esta fase, considerando el tiempo necesario para 
la investigación, la elaboración de documentos de diseño y la planificación detallada del proyecto.
Además, todo el trabajo previo hace posible esta aparentemente corta planificación, ya que se han realizado
tareas de investigación y recopilación de información durante el periodo previo al inicio oficial del proyecto.

\subsection{Fase 2: Diseño e Implementación}

Durante esta fase se llevará a cabo el diseño incremental y la implementación de los componentes 
fundamentales del sistema, es decir, el núcleo jugable del proyecto. La fase se ha descompuesto en 
iteraciones ligadas al backlog para que cada issue quede asociado a un objetivo claro y a una razón 
concreta de cierre.

El objetivo de esta fase es establecer una base sólida sobre la cual se puedan construir las funcionalidades
adicionales, diferenciando entre features del FDD y tareas concretas del backlog. Así, por ejemplo, la feature
\gls{generacionprocedural} se materializa en dos tareas distintas, \textit{Generación Procedural (I)} y
\textit{Generación Procedural (II)}, porque la primera establece el esqueleto del sistema y la segunda lo
ajusta una vez integrados el jugador, los enemigos y la estructura de salas.

La implementación del núcleo se llevará a cabo de forma conjunta por los dos integrantes del equipo,
de forma que se facilite el trabajo y la colaboración. Las tareas de la fase se agrupan en seis iteraciones,
cada una con una justificación funcional propia:

\begin{enumerate}
\item \textbf{Iteración 1.} \textit{Generación Procedural (I)}, \textit{Añadir Jugador} y \textit{Añadir Enemigo}. Esta iteración fijó la base jugable mínima, porque sin generación de niveles y entidades principales no tenía sentido cerrar el resto del flujo de juego.
\item \textbf{Iteración 2.} \textit{Generación Procedural (II)}, \textit{Lógica de Círculos (I)} y \textit{Cámara del Jugador}. La segunda iteración se reservó para estabilizar la estructura del nivel antes de seguir con la navegación y el encuadre de la cámara.
\item \textbf{Iteración 3.} \textit{Lógica de Círculos (II)}, \textit{Dualidad de Personajes} y \textit{Minimap (I)}. Se pospuso parte de esta iteración porque la lógica de círculos necesitaba validarse tras la segunda pasada sobre la generación y porque el minimapa dependía de una estructura de nivel estable.
\item \textbf{Iteración 4.} \textit{Head-Up Display (HUD)}, \textit{Menús (I)} y \textit{Menús (II)}. La interfaz se abordó cuando ya existía una experiencia jugable básica, para no bloquear el trabajo de sistema con elementos puramente de presentación.
\item \textbf{Iteración 5.} \textit{Actualizar Sprites (I)}, \textit{Sistema de Combate (I)} y \textit{Sistema de Enemigos (I)}. El combate y el comportamiento enemigo se cerraron más tarde porque dependían de la cámara, la navegación y la base de interacción ya estabilizadas.
\item \textbf{Iteración 6.} \textit{Gestión de Inventario}, \textit{Objetos} y \textit{Minimap (II)}. Esta iteración se dejó para el final de la fase porque inventario y objetos solo tenían sentido con el resto del bucle jugable consolidado, y el minimapa necesitaba un repaso final tras integrar todos los sistemas previos.
\end{enumerate}

Se ha estimado una duración amplia para esta fase, considerando la complejidad de los componentes a implementar
y la necesidad de ir cerrando cada issue cuando realmente dejaba de depender de otra iteración anterior. La idea es
que ambos integrantes del equipo puedan colaborar y compartir el conocimiento, facilitando el trabajo de la fase de cierre.

\subsection{Fase 3: Cierre del Proyecto}

En esta fase se reservará el espacio para la revisión final del proyecto, la preparación de la defensa,
la consolidación de la documentación y el cierre administrativo de la memoria. De momento se deja vacía,
porque el trabajo asociado a este tramo todavía no se ha incorporado al texto final.

Se ha estimado una duración de dos semanas para esta fase, considerando que la memoria se encontrará en
un estado avanzado al trabajar sobre ella a lo largo de todo el proyecto, y que la preparación de la
presentación final requerirá tiempo para organizar el contenido y ensayar la exposición.

\section{Informe de tiempos del proyecto}

Lo mismo que el anterior pero con datos reales. Ver Tabla \ref{tab:InformeTiempos}.

\begin{table*}[htb]
	\centering
	\begin{coolTable}{ll}{2}
{Resumen de iteraciones}
	\textbf{Iteración 1}&10/10/14 a 21/10/14\\
	\textbf{Iteración 2}&21/10/14 a 15/11/14\\
	\textbf{...}&dd/mm/aa a dd/mm/aa\\
	\end{coolTable}
	\caption{Planificación temporal de iteraciones\label{tab:InformeTiempos}}
\end{table*}

Justificar los retrasos de forma detallada aquí para cada una de las iteraciones. Explicar las razones.

\section{Desglose por fases}

En esta sección se recogerá el contenido que se vaya completando para cada fase del proyecto. 
El objetivo es dejar un espacio claro en la memoria para ir incorporando, de forma ordenada, 
qué se ha hecho en cada momento, qué decisiones se han tomado y qué resultados se han obtenido.

\subsection{Fase 1: Planificación y Gestión}
En esta subsección se describirá todo lo relacionado con el arranque del proyecto: análisis inicial, 
definición de objetivos, elección de herramientas, organización del backlog y preparación del plan de trabajo.

\subsection{Fase 2: Diseño e Implementación}
En esta subsección se recogerá el desarrollo de las mecánicas base del juego, distribuidas por iteraciones,
junto con la justificación de por qué cada issue se cerró en una iteración concreta y no en la anterior.

\subsection{Fase 3: Cierre del Proyecto}
En esta subsección se añadirá, cuando corresponda, la documentación final, la preparación de la presentación 
y la revisión final del trabajo antes de su entrega.